% !PRINTED 246
\weberpage{246}
%<<CONT
im gewöhnlichen Sinne verstehen. Alle Zahlen, auf die man hierbei kommt, lassen sich auf die Form

\weq{2}{\alpha = a_0 + a_1\varepsilon + \cdots + a_{n-1}\varepsilon^{n-1}}

bringen. Da bei der Rechnung mit solchen Zahlen der Modul $p$ immer festgehalten wird, so wollen wir zwei Zahlen, die nach dem Modul $p$ congruent sind, geradezu als \gesperrt{gleich} bezeichnen. Objecte der Rechnung sind dann eigentlich nicht die einzelnen Zahlen selbst, sondern die aus allen unter einander congruenten Zahlen bestehenden \gesperrt{Zahlclassen}. Diese Zahlen $\alpha$ werden die \gesperrt{Galois'schen Imaginären} genannt. Das zu einem bestimmten $\varepsilon$ gehörige System solcher imaginärer Zahlen bezeichnen wir der Abkürzung wegen mit $\mathfrak{C}$.

Wir haben dann zunächst den Satz:

\gesperrt{2. Es giebt $p^n$ und nicht mehr verschiedene Zahlen in $\mathfrak{C}$.}

Aus dem Satze 1. aber folgt noch:

\gesperrt{3. Das Product von zwei oder mehr Zahlen aus $\mathfrak{C}$ ist dann und nur dann gleich Null, wenn wenigstens einer der Factoren gleich Null ist.}

Wir erhalten das vollständige Zahlensystem $\mathfrak{C}$, wenn man die rationalen Zahlen $a_0, a_1, \ldots, a_{n-1}$ in (2) je ein volles Restsystem nach dem Modul $p$ durchlaufen lässt. Jedes System $\mathfrak{C}$ enthält die Reste der natürlichen Zahlen $0, 1, \ldots, p-1$, und für $n = 1$ ist $\mathfrak{C}$ mit diesem Systeme identisch.

Ist $\alpha$ eine feste von Null verschiedene Zahl in $\mathfrak{C}$, so durchläuft $\alpha\xi$ zugleich mit $\xi$ das volle System $\mathfrak{C}$. Denn aus 3. folgt, dass $\alpha\xi$ nur dann gleich $\alpha\xi'$ sein kann, wenn $\xi = \xi'$ ist. Daraus folgt:

\gesperrt{4. Sind $\alpha, \beta$ zwei gegebene Zahlen in $\mathfrak{C}$ und $\alpha$ von Null verschieden, so giebt es eine und nur eine Zahl $\gamma$ in $\mathfrak{C}$, die der Bedingung}

\[
\alpha\gamma = \beta
\]

\gesperrt{genügt.}

Damit ist auch die Operation der Division in dem Systeme $\mathfrak{C}$ als erlaubt nachgewiesen. Wir bezeichnen die Zahl $\gamma$, deren Existenz in 4. ausgesprochen ist, mit

\[
\beta : \alpha \quad \text{oder} \quad \frac{\beta}{\alpha}.
\]
